\title{Direct Preference Optimization: Your Language Model is Secretly a Reward Model}

\begin{abstract}
Although large-scale unsupervised language models (LMs) acquire extensive world knowledge and some reasoning abilities, exerting precise influence over their outputs remains challenging, primarily because of the fully unsupervised framework used during their training. Current approaches that aim to enhance this controllability typically involve gathering human judgments that rank the quality of model outputs and subsequently adjusting the LM to reflect these preferences—most often through reinforcement learning from human feedback (RLHF). Nonetheless, the RLHF methodology is often intricate and unstable, as it requires the initial construction of a reward model that captures human preferences, followed by reinforcement learning-based fine-tuning of the LM to optimize this reward, all while avoiding significant divergence from the original model. In this work, we propose an alternative reward model parameterization for RLHF that allows the derivation of the associated optimal policy in a closed-form solution. This formulation makes it possible to address the conventional RLHF objective via a straightforward classification loss. We term our method \textit{Direct Preference Optimization} (DPO), and it is distinguished by its robustness, high performance, and computational efficiency. DPO removes the necessity for sampling from the LM during fine-tuning and dispenses with extensive hyperparameter searches. Experimental results indicate that DPO matches or surpasses the effectiveness of prevailing techniques for aligning LMs with human preferences. Remarkably, DPO outperforms RLHF using PPO in steering sentiment in generated text, and achieves equivalent or superior results in summarization and single-turn dialogue tasks, all while offering significantly greater simplicity in both implementation and training.
\end{abstract}

\section{Introduction}
Large-scale unsupervised language models (LMs) trained on extensive corpora demonstrate a range of remarkable capabilities~\citep{chowdhery2022palm, brown2020language, touvron2023llama,bubeck2023sparks}. Nonetheless, these models are exposed to data authored by humans with differing intentions, priorities, and competencies. Consequently, not all traits reflected in the data are desirable for the model to emulate; for example, although we may prefer our AI coding assistant to \textit{recognize} prevalent programming errors in order to suggest corrections, we ultimately wish for the model's code outputs to reflect the (potentially infrequent) instances of superior coding practices found in the training set. Analogously, while it is advantageous for a language model to be \textit{informed} about a widespread misconception held by half of the population, it is undesirable for the model to reproduce this false belief in half its responses to related queries. Thus, carefully identifying and promoting the model’s \emph{targeted responses and conduct}—as opposed to the entire spectrum of its \textit{acquired knowledge and capabilities}—is fundamental to creating AI systems that are secure, effective, and manageable \citep{ouyang2022training}. Although popular approaches leverage reinforcement learning (RL) to align LMs with human preferences, we demonstrate that the objective optimized in existing RL-based methods can be solved exactly via a straightforward binary cross-entropy loss, significantly streamlining the process of learning from preferences.

In general, prior strategies introduce desired attributes into a language model by employing carefully constructed collections of human preferences, which articulate the behaviors considered beneficial or safe by people. This process of preference optimization follows an initial phase of extensive unsupervised pre-training on massive textual datasets. Although the most direct method for preference adaptation involves supervised fine-tuning using exemplary human-provided outputs, the most effective approaches have been reinforcement learning from human or AI feedback (RLHF/RLAIF; \citep{christiano2017deep,bai2022constitutional}). RLHF entails constructing a reward model from annotated human preferences and then applying RL to guide the language model's policy toward producing high-reward outputs without deviating significantly from its original parameters. Despite RLHF’s success in generating models with outstanding conversational and programming abilities, the RLHF workflow is substantially more intricate than that of standard supervised learning, requiring the training of several LMs and online sampling from the policy during optimization, which considerably increases computational overhead.

This work demonstrates an approach for optimizing language models to align with human preferences directly, bypassing the need for explicit reward modeling or reinforcement learning frameworks. We introduce \textit{{Direct Preference Optimization} (DPO)}, a method that implicitly targets the same objective as standard RLHF approaches—maximizing reward subject to a KL-divergence constraint—while remaining easy to implement and train. Conceptually, DPO adjusts the model by boosting the relative log likelihood of favored over unfavored responses, utilizing a per-instance importance weighting scheme that guards against the model collapse issues observed with straightforward probability ratio objectives. Like established approaches, DPO makes use of a formal preference model (e.g., the Bradley-Terry model; \cite{bradley1952rankanalysis}) to evaluate the concordance between reward functions and observed preference data. Unlike previous methods, which leverage the preference model first to establish a preference-based loss for reward model training and then optimize a policy to maximize this reward, DPO reformulates the problem: it expresses the preference loss directly as a function of the policy through a variable transformation. Consequently, given a corpus of human-annotated preferences among model outputs, DPO can be trained using a binary cross-entropy loss, thereby identifying the optimal policy corresponding to an implicit reward structure that matches the empirical preferences.

The primary contribution of this paper is Direct Preference Optimization (DPO), an RL-free and straightforward approach for preference-driven language model training. Our experimental results indicate that DPO achieves performance at least on par with current techniques—including RLHF approaches based on PPO—across applications like sentiment adjustment, summarization, and conversational tasks, tested on language models up to 6B parameters.

Large self-supervised language models have demonstrated the ability to tackle some tasks without prior examples \citep{radford2019language} or by using only a handful of examples provided as prompts \citep{gpt3,megatron,chowdhery2022palm}. Nonetheless, their effectiveness on various applications and their alignment with user objectives can be greatly enhanced through fine-tuning on instruction-based datasets accompanied by human-authored completions \citep{mishra-etal-2022-cross,sanh2022multitask,chung2022scaling,thoppilan2022lamda}. This approach, termed 'instruction-tuning,' facilitates LLMs’ capacity to extrapolate to novel instructions beyond those seen in the instruction-tuning set and generally boosts practical usefulness \citep{chung2022scaling}. Although instruction tuning is successful, gathering \textit{comparative} human assessments of completion quality is often more feasible than collecting expert-created demonstrations. As a result, recent studies have explored fine-tuning LLMs on datasets constructed from human preferences, leading to enhanced abilities in domains such as translation \citep{kreutzer-etal-2018-reliability}, summarization \citep{stiennon2022learning,ziegler2020finetuning}, narrative generation \citep{ziegler2020finetuning}, and following instructions \citep{ouyang2022training,ramamurthy2023is}. These approaches typically involve first training a neural reward model aligned with collected preferences, leveraging preference models like the Bradley-Terry model \citep{bradley1952rankanalysis}, and subsequently optimizing the language model with respect to this learned reward signal using reinforcement learning methods such as REINFORCE \citep{williams1992reinforce}, proximal policy optimization (PPO; \cite{schulman2017proximal}), or related algorithms \citep{ramamurthy2023is}. A closely affiliated research direction uses LLMs that have undergone instruction-following fine-tuning with human feedback to synthesize additional preference datasets that target attributes like safety or harmlessness \citep{bai2022constitutional}, relying on limited human supervision through written rubrics to guide model annotation. This body of work effectively merges two research streams: one on employing reinforcement learning to optimize language models for diverse goals~\citep{Ranzato2015SequenceLT,paulus2018a,wu2018learning}, and another on generic methodologies for learning from human preferences \citep{christiano2017deep,kupcsik2018learning}. Despite the attractiveness of relative preference learning, the process of reinforcement learning-based fine-tuning for large language models continues to pose substantial practical difficulties; the present work contributes a theoretically grounded solution for optimizing relative preference signals without requiring reinforcement learning.

Research into policy learning from preferences, outside the realm of language, has been conducted in both bandit problems and reinforcement learning, yielding multiple methods. The contextual dueling bandit (CDB; \cite{yue2012karmed,dudik2015contextual}) framework, for instance, investigates contextual bandit learning where the feedback consists of preferences or rankings over actions instead of direct rewards. When absolute reward values are unavailable, theoretical treatments of CDBs introduce the concept of a \textit{von Neumann winner}, defined as a policy whose expected likelihood of outperforming any other policy is at least 50\% \citep{dudik2015contextual}. Nevertheless, CDB approaches typically collect preference data in an online fashion, whereas preference learning from human annotators commonly proceeds from a static dataset containing pairs of actions labeled with preferences \citep{yan2022human}. In a similar vein, \textit{preference-based RL} (PbRL) uses binary preferences derived from an \textit{unknown} ‘scoring’ or reward function rather than relying on explicit rewards \citep{BusaFekete2014,ruiz2023dueling}. PbRL has seen a range of algorithms, including those allowing the re-utilization of previously gathered off-policy preference data. Most of these methods, however, involve an explicit estimation of the hidden scoring function (the reward model) as a separate stage prior to its optimization \citep{jain2013learning,BusaFekete2014,christiano2017deep,sadigh2017active,kupcsik2018learning}. Contrasting with these, our work introduces a unified approach that optimizes the policy directly to align with observed preferences, obviating the need for an intermediate reward modeling step.

\section{Preliminaries}\label{section:prelims}

Here, we summarize the RLHF process as described by \citeauthor{ziegler2020finetuning} and further elaborated in later works \citep{stiennon2022learning, bai2022training, ouyang2022training}. The procedure generally consists of three core steps: (1) supervised fine-tuning (SFT), (2) collection of preference samples and training a reward model, and (3) reinforcement learning optimization.

\textbf{SFT}: The initial step in RLHF typically involves adapting a pre-trained language model to downstream applications such as dialogue or summarization via supervised learning on curated datasets, resulting in a fine-tuned model $\pi^\text{SFT}$.

\textbf{Reward Modelling Phase}: Next, prompts $x$ are fed into the SFT model, which generates pairs of outputs $(y_1, y_2)\sim \pi^\text{SFT}(y \mid x)$. Human evaluators then provide pairwise judgments, indicating their preference for one output, which is recorded as $y_w\succ y_l \mid x$, with $y_w$ denoting the favored and $y_l$ the less favored response from the pair. The assumption is that these preferences are governed by some underlying, inaccessible reward function $r^*(y, x)$. To model these judgments, the Bradley-Terry (BT) \cite{bradley1952rankanalysis} model is frequently employed, though the framework can accommodate more general ranking models such as the Plackett-Luce \citep{plackett1975analysis, luce2012individual} when preferences over multiple alternatives are available. The BT formulation posits that the human preference probability distribution $p^*$ is described as follows:

\begin{equation}\label{eq:bradley-terry}
    p^*(y_1\succ y_2 \mid x)=\frac{\exp\left(r^*(x, y_1)\right)}{\exp\left(r^*(x, y_1)\right) + \exp\left(r^*(x, y_2)\right)}.
\end{equation}

Given a fixed dataset of pairwise preferences, $\mathcal{D}=\bigl\{x^{(i)}, y_w^{(i)}, y_l^{(i)}\bigr\}_{i=1}^N$, assumed to be drawn according to $p^*$, we introduce a parameterized reward model $r_{\phi}(x, y)$ and seek to estimate its parameters by maximizing the likelihood. This task can be cast as a binary classification, where the associated loss is the negative log-likelihood:

\begin{equation}\label{eq:reward_model}
    \mathcal{L}_R(r_{\phi}, \mathcal{D}) = -\mathbb{E}_{(x, y_w, y_l)\sim \mathcal{D}}\bigl[\log \sigma(r_{\phi}(x, y_w)- r_{\phi}(x, y_l))\bigr]
\end{equation}

Here, $\sigma$ denotes the logistic sigmoid. In large language model settings, $r_{\phi}(x, y)$ is typically initialized from a supervised fine-tuned model $\pi^\text{SFT}(y \mid x)$, with an additional linear projection atop the final transformer layer output, resulting in a scalar reward prediction \cite{ziegler2020finetuning}. To control for reward estimate variance, previous work standardizes rewards to satisfy $\mathbb{E}_{x,y\sim \mathcal{D}}\left[r_\phi(x, y)\right] = 0$ across $x$.

\textbf{RL Fine-Tuning Phase}: Within the reinforcement learning step, the trained reward model guides the language model's updates. Following earlier studies~\citep{jaques2017sequence, jaques2020human}, optimization is framed as

\begin{equation}\label{eq:RL}
\max_{\pi_{\theta}}  \mathbb{E}_{x\sim \mathcal{D}, y\sim \pi_{\theta}(y \mid x)}\bigl[r_{\phi}(x, y)\bigr] - \beta\mathbb{D}_{\textrm{KL}}\bigl[\pi_{\theta}(y\mid x)\mid \mid \pi_\text{ref}(y\mid x)\bigr],
\end{equation}

where the scalar $\beta$ balances proximity to a reference policy $\pi_\text{ref}$, commonly the pre-trained or SFT model $\pi^\text{SFT}$. Conventionally, the initial policy $\pi_\theta$ is also set to $\pi^\text{SFT}$. This regularization is essential, as it constrains the model to remain near distributions for which the reward model is reliable, preserves sample diversity, and mitigates collapse to a limited set of high-reward responses. Since the generative process is discrete, the objective cannot be directly differentiated and is handled through reinforcement learning techniques. The prevalent method \citep{ziegler2020finetuning, stiennon2022learning, bai2022training, ouyang2022training} forms a reward of the form ${r(x, y) = r_{\phi}(x, y) -\beta (\log \pi_{\theta}(y\mid x) - \log \pi_\text{ref}(y\mid x))}$, optimizing it via PPO \cite{schulman2017proximal}.

\section{Direct Preference Optimization}\label{sec:DPO}

Given the substantial complexities involved in applying reinforcement learning for fine-tuning large models, our objective is to introduce a streamlined alternative for optimizing policies directly from preference data. Contrary to standard RLHF pipelines, which involve explicit reward model training followed by reinforcement learning, we utilize a particular parameterization of the reward model that yields the optimal policy in closed-form—obviating the need for iterative RL. As elaborated below, our principal insight is to exploit a tractable relationship between rewards and the optimal policy, thereby converting a loss originally defined on rewards into an equivalent loss on the policy itself. This reparameterization eliminates the requirement for a separately trained reward model, yet remains consistent with commonly adopted human preference models, such as the Bradley-Terry framework. In summary, this approach allows the policy network to act as both the language model and the implicit reward estimator.

\textbf{DPO Objective Formulation.} To begin, we consider the standard RL objective as established in prior research, denoted in Eq.~\ref{eq:RL}, for a generic reward function $r$. Previous work~\citep{peters2007reinforcement, peng2019advantage, korbak2022reinforcement, go2023aligning} demonstrates that the solution maximizing the reward subject to a KL constraint (Eq.~\ref{eq:RL}) admits the following structure:

\begin{equation}\label{eq:op_policy}
    \pi_r(y\mid x) = \frac{1}{Z(x)}\pi_\text{ref}(y\mid x)\exp\left(\frac{1}{\beta}r(x, y)\right),
\end{equation}

where the partition function $Z(x) =\sum_{y}\pi_\text{ref}(y\mid x)\exp\left(\frac{1}{\beta}r(x, y)\right)$ ensures proper normalization. A complete derivation can be found in Appendix \ref{app:derivation1}. When substituting the MLE estimate $r_{\phi}$ for the true reward function $r^*$, direct evaluation of the partition function $Z(x)$ remains computationally demanding~\citep{korbak2022reinforcement, go2023aligning}, thereby impeding practical use of this expression. Nevertheless, Eq.~\ref{eq:op_policy} can be algebraically rearranged to solve for the reward function in terms of its associated optimal policy $\pi_r$, the reference distribution $\pi_\text{ref}$, and the undetermined normalization constant $Z(\cdot)$. Specifically, applying the logarithm to both sides of Eq.~\ref{eq:op_policy} and rearranging yields:

\begin{equation}\label{eq:main_eq}
    r(x,y) =\beta \log \frac{\pi_r(y\mid x)}{\pi_\text{ref}(y\mid x)} + \beta \log Z(x).
\end{equation}

This reparameterization is equally applicable to the true reward $r^*$ and its associated optimal policy $\pi^*$. The Bradley-Terry model conveniently depends only on differences between rewards assigned to two candidate responses, namely ${p^*(y_1 \succ y_2 \mid x) = \sigma(r^*(x, y_1) - r^*(x, y_2))}$. By plugging the reparameterized reward from Eq.~\ref{eq:main_eq} into the preference model Eq.~\ref{eq:bradley-terry}, the normalization term cancels out, allowing the probability of observed preferences to be stated exclusively using $\pi^*$ and $\pi_\text{ref}$. Consequently, under the Bradley-Terry framework, the optimal RLHF policy $\pi^*$ fulfills:

\begin{equation}\label{eq:objective}
    p^*(y_1\succ y_2 \mid x)=\frac{1}{1 + \exp\left(\beta \log \frac{\pi^*(y_2\mid x)}{\pi_\text{ref}(y_2\mid x)} - \beta \log \frac{\pi^*(y_1\mid x)}{\pi_\text{ref}(y_1\mid x)}\right)}
\end{equation}

Further details are available in Appendix~\ref{app:derivation2}. Although Eq.~\ref{eq:objective} is expressed in terms of the Bradley-Terry model, an analogous formulation is attainable for broader preference models such as Plackett-Luce~\citep{plackett1975analysis, luce2012individual}, as presented in Appendix~\ref{app:plackett_luce_models}.

With the likelihood of human-labeled preference data now depending on the optimal policy rather than the reward function, we can cast a maximum likelihood estimation problem for a parameterized policy $\pi_\theta$. Similar to the approach used in reward modeling (see Eq.~\ref{eq:reward_model}), the policy learning objective becomes:

\begin{equation}\label{eq:optimum_model}
    \mathcal{L}_\text{DPO}(\pi_{\theta}; \pi_\text{ref}) = -\mathbb{E}_{(x, y_w, y_l)\sim \mathcal{D}}\left[\log \sigma \left(\beta \log \frac{\pi_{\theta}(y_w\mid x)}{\pi_\text{ref}(y_w\mid x)} - \beta \log \frac{\pi_{\theta}(y_l\mid x)}{\pi_\text

By adopting this approach, we model an implicit reward through a different parameterization, such that the optimal policy corresponds directly to $\pi_\theta$. Furthermore, our methodology is tantamount to learning a reparameterized Bradley-Terry model, granting it desirable theoretical attributes—such as consistency given appropriate assumptions on the distribution of preference data \cite{bong2022generalized}. Additional theoretical analysis of DPO and its connections to related work can be found in Section~\ref{sec:theory}.

\textbf{Mechanistic perspective on the DPO update.} To gain deeper intuition about the DPO update, one may inspect the gradient of the loss $\mathcal{L}_\text{DPO}$. With respect to the parameter vector $\theta$, the gradient takes the following form:

\begin{multline*}\label{eq:gradient}
    \nabla_\theta \mathcal{L}_\text{DPO}(\pi_\theta;\pi_\text{ref}) = \\ -\beta\mathbb{E}_{(x, y_w, y_l) \sim \mathcal{D}} \bigg[\underbrace{\sigma(\hat{r}_\theta(x, y_l) - \hat{r}_\theta (x, y_w))}_\text{emphasizes mistakes of reward model}\bigg[\underbrace{\nabla_\theta\log \pi(y_w \mid x)}_\text{promote $y_w$} - \underbrace{\nabla_\theta\log\pi(y_l \mid x)}_\text{penalize $y_l$}\bigg]\bigg],
\end{multline*}

Here, $\hat{r}_\theta(x, y) = \beta \log \frac{\pi_\theta(y \mid x)}{\pi_\text{ref}(y \mid x)}$ defines an implicit reward ascribed by the policy $\pi_\theta$ relative to a reference model $\pi_\text{ref}$ (as further examined in Section~\ref{sec:theory}). At an intuitive level, the gradient of $\mathcal{L}_\text{DPO}$ acts to boost the probability of the favored completion $y_w$ while diminishing that of the unfavored completion $y_l$. Critically, the update is modulated according to the extent to which the implicit reward function $\hat{r}_\theta$ erroneously scores $y_l$ above $y_w$—this magnitude is scaled by $\beta$ and reflects both the misordering and the effect of the KL regularization. Our experimental findings indicate that this specific weighting is vital: omitting the weighting factor (i.e., applying the method na\"ively) often results in degeneration of the language model, as illustrated in Appendix Table~\ref{tab:unlikelihood_generations}.

\textbf{DPO overview.} The standard DPO procedure involves the following steps: 1) For each prompt $x$, generate sample completions $y_1, y_2 \sim \pi_\text{ref}(\cdot \mid x)$, assign human preference labels, and compile the offline preference dataset $\mathcal{D} = \{x^{(i)}, y_w^{(i)}, y_l)^{(i)}\}_{i=1}^N$; 2) Train the language model $\pi_\theta$ by minimizing $\mathcal{L}_\text{DPO}$, utilizing $\pi_\text{ref}$, $\mathcal{D}$, and a chosen $\beta$. 
Practically, there is a strong incentive to leverage existing publicly available preference datasets instead of generating new samples and annotations. As these datasets are collected with $\pi^\text{SFT}$, $\pi_\text{ref}$ is initialized as $\pi^\text{SFT}$ whenever such a model exists. In scenarios where $\pi^\text{SFT}$ is absent, $\pi_\text{ref}$ is initialized by maximizing the likelihood over preferred completions ${(x, y_w)}$, specifically, by setting ${\pi_\text{ref} = \argmax_{\pi}\mathbb{E}_{x, y_w \sim \mathcal{D}}\left[\log \pi(y_w \mid x)\right]}$. This approach helps to alleviate the potential mismatch between the actual reference distribution (which cannot be directly accessed) and the $\pi_\text{ref}$ employed by DPO. Additional implementation specifics and choices of hyperparameters are discussed in Appendix~\ref{app:implementation}.

\section{Theoretical Analysis of DPO} This section provides a deeper theoretical interpretation of the DPO approach, establishes foundational results, and connects the strengths of DPO with the common shortcomings found in actor-critic methods employed in RLHF (such as PPO~\cite{schulman2017proximal}).

\label{sec:theory}

\subsection{Your Language Model Is Secretly a Reward Model} DPO can circumvent both the need for fitting a separate reward model and running reinforcement learning to train the policy, by relying on a unified maximum likelihood training objective. Importantly, the optimization objective in Eq. \ref{eq:main_eq} aligns with the Bradley-Terry framework using a reward representation $r^*(x, y) = \beta \log\frac{\pi^*_\theta(y \mid x)}{\pi_\text{ref}(y \mid x)}$. Optimization over the parameterized model $\pi_{\theta}$ then coincides with optimizing a reward model as in Eq. \ref{eq:reward_model}, modulo a change of variables. Here, we develop the theoretical grounding for this reparameterization, demonstrate that it retains full expressivity over the space of possible reward models, and show that it facilitates the exact recovery of the optimal policy. We begin by introducing an equivalence relation for reward functions.

\begin{definition}
Two reward functions $r(x, y)$ and $r'(x, y)$ are said to be equivalent if and only if ${r(x, y)-r'(x, y) = f(x)}$ for some function $f$.     
\end{definition}

This definition clearly forms an equivalence relation, thus dividing the space of reward functions into distinct equivalence classes. With this, we can establish the following two lemmas:

\begin{lemma}\label{lemma:same_prefrence}
Within the Plackett-Luce preference framework, which includes the Bradley-Terry model as a special case, any two reward functions belonging to the same equivalence class produce identical preference distributions.
\end{lemma}

\begin{lemma}\label{lemma:same_policy}
    If two reward functions are in the same equivalence class, they will yield the same optimal policy when solving the constrained reinforcement learning problem.
\end{lemma}

The proofs of these lemmas are elementary and are provided in Appendix \ref{app:lemma1}. The first lemma highlights the well-established under-specification property inherent to Plackett-Luce-type models \cite{plackett1975analysis}. Owing to this non-identifiability, it is common to introduce additional constraints to make the maximum likelihood estimation in Eq. \ref{eq:reward_model} well-posed \cite{bong2022generalized}. The second lemma asserts that all reward functions within an equivalence class induce the same optimal policy; therefore, for our purposes, recovering any representative from the optimal equivalence class suffices. The following theorem, demonstrated in Appendix~\ref{app:thm1}, formalizes this insight:

\begin{theorem}\label{thm:main}
    Provided certain mild conditions, every reward class that is consistent with Plackett-Luce (including Bradley-Terry) models can be written in the form ${r(x, y) = \beta \log \frac{\pi(y\mid x)}{\pi_\text{ref}(y\mid x)}}$ for some model $\pi(y\mid x)$ and a specified reference model $\pi_\text{ref}(y \mid x)$.
\end{theorem}

\begin{sproof}
    Let $r(x, y)$ denote an arbitrary reward function, associated with an optimal model $\pi_r(y \mid x)$ defined by Eq. \ref{eq:op_policy}. We aim to demonstrate that there exists a reward function in the equivalence class of $r$ that admits the reparameterized form given above. Define the projection map $f$ as  
\begin{equation}
    f(r; \pi_\text{ref}, \beta)(x, y) = r(x, y) - \beta\log\sum_{y}\pi_\text{ref}(y\mid x)\exp\left(\frac{1}{\beta}r(x, y)\right)
\end{equation}
The function $f$ normalizes $r(x, y)$ by subtracting the log-partition function of $\pi_r$, which depends only on the input $x$. Thus, $f(r; \pi_\text{ref}, \beta)(x, y)$ lies within the equivalence class of $r(x, y)$. By substituting the right-hand side of Eq.~\ref{eq:main_eq} (applicable to any $r$), we see that $f(r; \pi_\text{ref}, \beta)(x, y) = \beta \log \frac{\pi_r(y\mid x)}{\pi_\text{ref}(y\mid x)}$. Consequently, this projection provides a canonical representative of $r$’s equivalence class with the sought-after form, and the reparameterization does not constrain the expressivity of the reward model.
\end{sproof}

An alternative perspective on Theorem~\ref{thm:main} is that it characterizes the precise reward function within each equivalence class that the DPO reparameterization picks, specifically, the reward function that fulfills:

\begin{equation}\label{eq:lag_p}
     \sum_{y}\underbrace{\pi_\text{ref}(y\mid x)\exp\left(\frac{1}{\beta}r(x, y)\right)}_{=\pi(y\mid x)\text{, using Thm.~\ref{thm:main} reparam.}} = 1,
\end{equation}

which ensures that $\pi(y\mid x)$ forms a legitimate probability distribution (namely, all probabilities are non-negative and total to 1). Moreover, referring to Eq.~\ref{eq:op_policy}, it becomes evident that Eq.~\ref{eq:lag_p} corresponds to the partition function of the optimal policy dictated by the reward function $r(x, y)$. A central realization behind the DPO algorithm is that by enforcing particular constraints on the otherwise underdetermined Plackett-Luce (and specifically, Bradley-Terry) family of preference models, we retain the same expressive power for reward modeling, while also ensuring that the optimal policy given by Eq.~\ref{eq:op_policy} is tractable in closed-form for any prompt $x$.

\subsection{Instability of Actor-Critic Algorithms} 

Our framework further serves as a lens through which we can understand instabilities associated with popular actor-critic methods, such as PPO, as used in RLHF. Adhering to the RLHF protocol, we concentrate on the RL fine-tuning phase as discussed in Section \ref{section:prelims}. By leveraging the connection to the control as inference paradigm \cite{levine2018reinforcement}, which applies to the constrained RL problem stated in \ref{eq:RL}, we consider a model parameterized by $\pi_{\theta}(y\mid x)$ and seek to minimize the divergence $\mathbb{D}_{\text{KL}}[\pi_{\theta}(y|x) \mid \mid \pi^*(y\mid x)]$, where $\pi^*$ represents the optimal policy described in Eq.~\ref{eq:optimum_model} and is determined by the reward $r_{\phi}(y, x)$. Through straightforward manipulations, this objective yields the following optimization formulation:

\begin{equation}\label{eq:AC}
    \max_{\pi_{\theta}}\mathbb{E}_{\pi_{\theta}(y\mid x)}\bigg[\underbrace{r_{\phi}(x, y) -\beta\log\sum_{y}\pi_\text{ref}(y\mid x)\exp\left(\frac{1}{\beta}r_{\phi}(x, y)\right)}_{f(r_{\phi}, \pi_\text{ref}, \beta)} - \underbrace{\beta\log\frac{\pi_{\theta}(y\mid x)}{\pi_\text{ref}(y\mid x)}}_{\text{KL}}\bigg]
\end{equation}

This objective has been the focus of optimization in previous studies \citep{ziegler2020finetuning, stiennon2022learning, bai2022training, ouyang2022training}, where the DPO-equivalent reward within the reward class $r_{\phi}$ is employed. Within this context, the normalization component in $f(r_{\phi}, \pi_\text{ref}, \beta)$ can be viewed as the soft value function corresponding to the reference policy $\pi_\text{ref}$. Although the normalization term does not change the ultimate optimal solution, omitting it may result in a policy gradient with considerable variance, thereby destabilizing the learning process. This normalization can be accounted for through a learned value function; however, such optimization presents its own challenges. As an alternative, earlier works have addressed normalization by using human completion baselines, effectively leveraging a single sample as a Monte Carlo approximation of the normalization term. Notably, the DPO reparameterization circumvents the need for any such baselines by yielding a reward function inherently normalized.

\section{Experiments}
In what follows, we empirically assess the effectiveness of DPO in learning policies from preference data alone. To begin, in a controlled text-generation environment, we investigate the trade-off efficiency between reward maximization and KL-divergence minimization with respect to the reference policy, comparing DPO to standard preference-based methods like PPO. Subsequently, we test DPO on more demanding settings involving larger model architectures and challenging RLHF tasks such as dialogue and summarization. Our findings indicate that, with minimal hyperparameter adjustment, DPO consistently matches or outperforms competitive approaches—including RLHF with PPO and the strategy of returning the top-performing trajectory among $N$ candidates evaluated under a learned reward model. Prior to discussing the results, we provide an overview of our experimental procedures; further methodological specifics can be found in Appendix~\ref{app:exp_details}.

\textbf{Tasks.} Our study investigates three distinct open-domain text generation challenges. Across all tasks, models are trained to learn a policy from a set of preference data $\mathcal{D}=\bigl\{x^{(i)}, y_w^{(i)}, y_l^{(i)}\bigr\}_{i=1}^N$. For \textbf{controlled sentiment generation}, $x$ serves as the opening segment of a movie review taken from the IMDb dataset \cite{maas-EtAl:2011:ACL-HLT2011}, and the model's objective is to generate a continuation $y$ exhibiting positive sentiment. To facilitate a controlled experiment, we create preference pairs for generations by leveraging a pre-trained sentiment classifier, ensuring that $p(\text{positive}\mid x,y_w)>p(\text{positive}\mid x,y_l)$. For the SFT approach, GPT-2-large is fine-tuned until it converges on reviews sourced from the IMDb training split (see App~\ref{app:sentiment_details} for more specifics). In the case of \textbf{summarization}, $x$ refers to a Reddit forum post, and the policy is required to produce a summary $y$ capturing the key arguments in the post. Consistent with previous methodologies, we use the Reddit TL;DR summarization dataset \citep{volske-etal-2017-tl} together with a human preference dataset constructed by \citeauthor{stiennon2022learning}. Our SFT baseline is trained on human-generated summaries of forum posts,\footnote{\url{https://huggingface.co/CarperAI/openai_summarize_tldr_sft}} employing the TRLX library \citep{leandro_von_werra_2023_7790115} for RLHF. The preference data was curated by \citeauthor{stiennon2022learning}, drawn from outputs of a similar, though independently trained, SFT model. Lastly, for \textbf{single-turn dialogue}, $x$ represents a user's input—ranging from scientific queries to personal advice. Here, the task for the policy is to deliver an informative and engaging reply $y$; we use the Anthropic Helpful and Harmless dialogue dataset \citep{bai2022training}, which includes 170k conversations between humans and an AI assistant. Each exchange concludes with two model-generated responses, each labeled according to human preference. Since no dedicated SFT model exists for this data, we construct the SFT model by fine-tuning a publicly available language model exclusively on preferred completions.

\textbf{Evaluation.} Our evaluation framework employs two complementary methodologies. To assess the capability of each algorithm in achieving the constrained reward maximization goal, we utilize a controlled sentiment generation scenario, in which we plot each algorithm’s frontier in terms of realized reward versus KL-divergence from the reference policy. This is feasible because the ground-truth reward function, supplied by a sentiment classifier, is directly accessible. Conversely, in real-world contexts, the ground-truth reward is unavailable. Thus, we turn to the notion of \textit{win rate} against a baseline policy, enlisting GPT-4 as an automatic judge—serving as a stand-in for human raters—to evaluate the quality of summaries and the helpfulness of single-turn dialogue responses. For summarization tasks, the comparison is against the test set’s reference summaries; for dialogue, the baseline is the preferred response identified in the test dataset. Although prior work has established that language models can often function as superior automated evaluators compared to traditional metrics \citep{Chen2023ExploringTU}, we reinforce our reliance on GPT-4 through a human study described in Sec.~\ref{sec:human-judgments}. The results reveal a strong concordance between GPT-4 and human evaluations, frequently showing human alignment with GPT-4 that matches or even surpasses inter-annotator consistency among humans.

\textbf{Methods.} Beyond DPO, our investigation encompasses several established strategies for training language models aligned with human preferences. We apply zero-shot prompting with \textbf{GPT-J} \citep{gpt-j} for summarization, and employ 2-shot prompting with \textbf{Pythia-2.8B} \citep{biderman2023pythia} in the dialogue setting. Additionally, we assess the performance of the standard \textbf{SFT} model and the \textbf{Preferred-FT} variant, which is fine-tuned via supervised learning to generate the preferred completion $y_w$—using either the SFT output (for sentiment control and summarization) or a generic language model (for dialogue). Another related supervised approach is the \textbf{Unlikelihood} method~\citep{welleck2019neural}, which encourages the model to assign high probability to $y_w$ while actively suppressing the likelihood of $y_l$; an optional hyperparameter $\alpha\in[0,1]$ controls the strength of the `unlikelihood' term. We also implement \textbf{PPO} \citep{schulman2017proximal} using a reward model trained on preference data, and include \textbf{PPO-GT}, a variant that operates as an oracle by utilizing the ground-truth reward function available in controlled sentiment settings. For sentiment experiments, two PPO-GT variants are considered: a standard version \cite{leandro_von_werra_2023_7790115}, and a customized one that incorporates reward normalization and refined hyperparameter selection to maximize performance (these adjustments are likewise applied to the regular PPO when using learned rewards). Lastly, we include a \textbf{Best of $N$} baseline, wherein $N$ candidate responses are sampled from the SFT model (or from Preferred-FT in the dialogue task), and the response that maximizes the reward model trained on preference data is selected. While this method can yield strong results by separating reward model quality from PPO’s optimization process, its computational overhead is substantial, even for modest values of $N$, since every query at inference time necessitates sampling $N$ separate completions.

\subsection{How well can DPO optimize the RLHF objective?}

The standard RLHF approach utilizes a KL-constrained reward maximization objective that mediates between maximizing reward and maintaining proximity to a reference policy. Consequently, algorithm comparisons must consider both the achieved reward and the corresponding KL divergence; a modest increase in reward accompanied by a substantial increase in KL divergence may not be advantageous. In Figure~\ref{fig:frontier-tldr-main}, we present the tradeoff curve (reward versus KL) for several algorithms applied to the sentiment analysis scenario. For each method, we conduct several independent training sessions, each time adjusting a policy conservativeness hyperparameter (target KL $\in\{3,6,9,12\}$ for PPO, $\beta \in \{0.05,0.1,1,5\}$, $\alpha\in\{0.05,0.1,0.5,1\}$ for unlikelihood, and different random seeds for preferred-FT). In total, the sweep encompasses 22 distinct runs. At intervals of every 100 training steps, up to convergence, we assess each policy on a fixed test prompt set by measuring the mean reward (using the true reward function) and the mean sequence-level KL\footnote{Defined as the sum over the per-timestep KL divergences.} relative to the reference policy $\text{KL}\left(\pi\mid \mid \pi_\text{ref}\right)$. Our results demonstrate that DPO establishes a reward-KL frontier that is substantially more efficient than other methods, attaining superior reward values while maintaining lower KL divergence. This outcome is noteworthy for several reasons: most significantly, although DPO and PPO target an identical optimization objective, DPO consistently realizes a better balance of reward and KL (i.e., DPO strictly dominates PPO on the reward/KL tradeoff). Furthermore, DPO surpasses PPO’s performance even under conditions where PPO is supplied with ground-truth rewards (PPO-GT).

\subsection{Can DPO scale to real preference datasets?}
\label{sec:dpo-real-datasets}
We proceed by examining how well DPO performs when fine-tuned on tasks in summarization and single-turn dialogue. In the context of summarization, it has been observed that automatic measures like ROUGE often show limited alignment with human judgments~\citep{stiennon2022learning}, and previous research indicates that language models fine-tuned with PPO based on human feedback tend to produce higher-quality summaries. Our comparison spans various approaches by generating samples from the TL;DR summarization dataset’s test split and computing each method’s average win rate against reference completions within this set. We explore sampling completions using a range of temperatures from 0.0 to 1.0, reporting the win rates in Figure~\ref{fig:frontier-tldr-main} (right). For consistency, DPO, PPO, and Preferred-FT all fine-tune the same underlying GPT-J SFT model\footnote{\url{https://huggingface.co/CarperAI/openai_summarize_tldr_sft}}. The results indicate that DPO achieves a win rate near 61\% at a temperature of 0.0, surpassing PPO, which attains approximately 57\% under its best sampling temperature. Additionally, DPO reaches a higher maximum win rate than the $N$ baseline. We emphasize that DPO’s $\beta$ parameter was not extensively tuned, suggesting the current outcomes might not fully reflect DPO’s capability. Another notable observation is that DPO exhibits greater stability across varying sampling temperatures compared to PPO, whose performance at higher temperatures can decline to that of the unmodified GPT-J model. Meanwhile, Preferred-FT provides minimal improvements over SFT. We further perform a direct human evaluation comparing DPO and PPO in Section~\ref{sec:human-judgments}, finding that DPO samples at temperature 0.25 were chosen over PPO’s at the same temperature in 58\% of pairwise preferences.

For single-turn dialogues, our evaluation focuses on the portion of the test split in the Anthropic HH dataset \citep{bai2022training} where only one round of human-assistant exchange occurs. In assessing performance using GPT-4, we measure win rates for each approach against the reference completions deemed preferable in the test set. Given the lack of a standard SFT model for this benchmark, we initialize from a pre-trained Pythia-2.8B, apply Preferred-FT on the selected completions to produce an in-distribution reference model, and subsequently employ DPO for further training. Additionally, we benchmark against the strongest result among 128 Preferred-FT completions (noting that the Best of $N$ baseline's performance plateaus at $N = 128$, as detailed in Appendix Figure~\ref{fig:best-of-n}), and evaluate a 2-shot prompted Pythia-2.8B base model. Across all temperature settings for these methods, DPO consistently achieves comparable or superior outcomes. We also assess an RLHF model trained via PPO on the Anthropic HH dataset \footnote{\url{https://huggingface.co/reciprocate/ppo_hh_pythia-6B}}, provided by a reputable source \footnote{\url{https://github.com/CarperAI/trlx/tree/main/examples/hh}}. However, despite our efforts, no prompt or sampling temperature with PPO yields results that surpass the unmodified Pythia-2.8B model. Given our TL;DR results and recognizing that both Best of 128 and PPO target the same reward, we interpret the Best of 128 metric as a reasonable stand-in for PPO-level effectiveness. In summary, DPO emerges as the only method that is both computationally efficient and capable of surpassing the performance of preferred completions on the Anthropic HH dataset, matching or exceeding the results of the resource-intensive Best of 128 baseline. As further illustrated in Figure~\ref{fig:dialogue-main}, DPO also achieves its optimal performance in relatively few training iterations.

\subsection{Generalization to a new input distribution}

To extend the analysis of how PPO and DPO handle distributional changes, we assess the performance of the policies trained during the Reddit TL;DR summarization experiment on a different type of data—namely, news articles from the CNN/DailyMail test set \citep{nallapati-etal-2016-abstractive}. The same optimal sampling temperatures identified for TL;DR (0 and 0.25) are employed. The outcomes are summarized in Table~\ref{tab:ood}. As before, we determine the win rate of GPT-4 when compared with reference summaries, applying the same GPT-4 (C) prompt utilized for Reddit TL;DR, with “forum post” replaced by “news article.” On this distinct input distribution, DPO maintains a notable advantage over the PPO policy. These results provide preliminary support that DPO policies are at least as capable as PPO policies in transferring to new data domains, even in the absence of extra unlabeled Reddit TL;DR prompts that are available to PPO.

\subsection{Validating GPT-4 judgments with human judgments}
\label{sec:human-judgments}
A human subject evaluation is conducted to assess the trustworthiness of GPT-4's scoring, focusing on the TL;DR summarization task and utilizing two variants of GPT-4 prompts. The \textbf{GPT-4 (S)} (simple) prompt requests the annotator to judge which summary best captures the core information in the source post. In contrast, the \textbf{GPT-4 (C)} (concise) prompt also considers conciseness as a criterion, motivated by our observation that, under the \textbf{GPT-4 (S)} prompt, GPT-4 tends to favor longer and more repetitive summaries than do human raters. The full prompt text can be found in Appendix~\ref{app:prompts}. To ensure broad coverage of output quality, three methods are compared: the top-performing approach (DPO, temperature 0.25), the lowest-performing (PPO, temperature 1.0), and an intermediate method (SFT, temperature 0.25). Each is benchmarked against PPO with greedy sampling at its best temperature setting. Our analysis indicates that GPT-4's preferences, regardless of prompt used, generally align with those of humans to the same extent as inter-human agreement, implying GPT-4 can be considered a credible surrogate for human judgment (note that due to constraints on human raters, we acquire multiple human annotations only for DPO and PPO-1 comparisons). In aggregate, results produced using the \textbf{GPT-4 (C)} prompt correspond more closely to human judgment, and this prompt is therefore adopted for primary findings in Section~\ref{sec:dpo-real-datasets}. Further methodological details, including the web-based rating interface and the roster of human participants, are detailed in Appendix~\ref{app:human-study}.

\section{Discussion}
Preference-based learning presents a robust and scalable approach for developing effective and aligned language models. In this work, we proposed DPO—a straightforward methodology for training language models using preferences, sidestepping reinforcement learning entirely. Unlike conventional methods that reformulate the preference learning challenge as an RL problem to leverage pre-existing RL algorithms, DPO establishes a correspondence between language model policies and reward functions. This allows direct optimization of language models to conform to human preferences via a basic cross-entropy loss, dispensing with the complexities of reinforcement learning while maintaining expressiveness. DPO achieves comparable or superior results to leading RLHF approaches, such as those using PPO, and accomplishes this with almost no hyperparameter adjustment, thereby significantly lowering the entry barrier for preference-based language model training.

\textbf{Limitations \& Future Work.} This study highlights several directions for subsequent investigation. One area concerns the extent to which the DPO-trained policy generalizes to out-of-distribution data compared to policies derived from explicit reward learning; while preliminary findings indicate that DPO's generalization is on par with PPO-based methods, more systematic evaluation is required. Another question is whether methods like self-labeling using the DPO policy could harness unlabeled prompts as effectively as other strategies. Further, it remains to be seen how reward over-optimization may emerge in direct preference optimization settings, and whether the minor reduction in performance observed in Figure~\ref{fig:dialogue-main}-right exemplifies this issue. In addition, although our experiments encompass models with up to 6B parameters, scaling DPO to much larger, state-of-the-art models offers promising avenues for research. On the evaluation side, our analysis indicates that GPT-4-generated win rates are sensitive to prompt phrasing; future work could focus on devising optimal techniques for obtaining reliable judgments from automated evaluators. Lastly, there is considerable potential for extending DPO beyond the scope of language model training from human preferences to the broader domain of generative modeling across other modalities.